\label{sec:methods}

\subsection{CrossPoll Pipeline Overview}

CrossPoll consists of three stages: (1) \textit{joint pretraining} on multiple source crops, (2) \textit{targeted fine-tuning} on held-out crops with varying label budgets, and (3) \textit{evaluation} on test crops. Figure 1 illustrates the pipeline.

\begin{figure}[H]
  \centering
  \fbox{\includegraphics[width=0.95\textwidth]{../figures/pipeline.png}}
  \caption{\textbf{CrossPoll Pipeline}: (a) Joint pretraining on three source crops (Apple, Strawberry, Tomato); (b) Fine-tuning with varying label budgets (25--500) on held-out Kiwi crop; (c) Evaluation and comparison against baselines. The key insight is that features learned across multiple flower morphologies generalize better to new crops than COCO-pretrained models.}
  \label{fig:pipeline}
\end{figure}

\subsection{Stage 1: Joint Pretraining on Source Crops}

\subsubsection{Dataset Preparation}

We use three diverse source crops to pretrain the detector:

\begin{table}[H]
\centering
\caption{Source Crop Dataset Statistics}
\label{tab:source_crops}
\begin{tabular}{lccccc}
\toprule
\textbf{Crop} & \textbf{Flowers} & \textbf{Train Imgs} & \textbf{Validation Imgs} & \textbf{Color} & \textbf{Morphology} \\
\midrule
Apple & 1,127 & 3,500 & 2,000 & White/Pink & Clustered, delicate \\
Strawberry & 5,419 & 3,500 & 2,000 & White/Yellow & Dense, small \\
Tomato & 1,891 & 2,500 & 1,500 & Yellow & Delicate, sparse \\
\midrule
\textbf{Total} & \textbf{8,437} & \textbf{9,500} & \textbf{5,500} & & \\
\bottomrule
\end{tabular}
\end{table}

Images are annotated with bounding boxes around individual flowers. We adopt the YOLO format (class ID, normalized $x$, $y$, $w$, $h$ per image). Each source crop exhibits distinct visual characteristics:

\begin{itemize}
    \item \textbf{Apple}: Large white-pink flowers, clustered on branch nodes
    \item \textbf{Strawberry}: Small white flowers with yellow stamens, densely packed on runners
    \item \textbf{Tomato}: Yellow flowers, sparse distribution, delicate petals
\end{itemize}

These morphological differences are crucial: flowers vary in size (10--100 pixels), color palette, spatial density, and occlusion patterns. Training on all three simultaneously exposes the model to a diverse ``flower manifold'' that should improve generalization to the target (Kiwi) crop.

\subsubsection{Joint Training Procedure}

We jointly train YOLOv8-nano on all three source crops for 300 epochs. Training uses:

\begin{itemize}
    \item \textbf{Model}: YOLOv8-nano (3.01M parameters, 8.2 GFLOPs)
    \item \textbf{Batch size}: $-1$ (auto-computed per hardware; empirically 16 on Apple Silicon)
    \item \textbf{Optimizer}: AdamW with $\text{lr}_0 = 0.001$
    \item \textbf{Augmentation}: Mosaic (1.0), flip (0.5), HSV (0.015, 0.7, 0.4), perspective (0.0)
    \item \textbf{Validation interval}: Every 5 epochs on pooled validation set
    \item \textbf{Stopping criterion}: Patience = 20 (stop if validation mAP@50 doesn't improve for 20 epochs)
\end{itemize}

\textbf{Data mixing strategy}: Crops are randomly shuffled per batch (not stratified), allowing the model to learn crop-invariant flower features. Unlike crop-specific training, this sampling strategy forces the detector to learn features robust to visual variation.

\textbf{Motivation}: Joint training differs from sequential transfer (pretrain on Apple, transfer to Strawberry, transfer to Tomato). We hypothesize simultaneous training on diverse flowers provides broader feature learning than sequential transfer-to-transfer pipelines. Section~\ref{sec:results} validates this.

\subsection{Stage 2: Fine-tuning on Target Crop with Label Budgets}

\subsubsection{Target Crop: Kiwi}

Kiwi flowers are selected as the held-out test crop. Characteristics:

\begin{table}[H]
\centering
\caption{Hold-out Target Crop: Kiwi}
\label{tab:target_crop}
\begin{tabular}{lcccc}
\toprule
\textbf{Property} & \textbf{Value} & \textbf{Annotation} & \textbf{Train} & \textbf{Test} \\
\midrule
Total Images & 302 & -- & 181 & 121 \\
Total Flowers & 1,247 & -- & 755 & 492 \\
Color & White/Pale green & Distinct from sources & -- & -- \\
Morphology & Delicate, recurved sepals & Unique within dataset & -- & -- \\
Avg. Flower Size & 40 pixels & Smaller than Apple & -- & -- \\
\bottomrule
\end{tabular}
\end{table}

Kiwi flowers are visually distinct from source crops: pale white-green color, more delicate morphology with curved sepals, and slightly smaller size. This creates a realistic domain shift scenario.

\subsubsection{Fine-tuning Protocol}

For each label budget $B \in \{25, 50, 100, 200, 500\}$:

\begin{enumerate}
    \item \textbf{Subsample training}: Randomly select $B$ images from Kiwi train set (181 available). When $B \geq 181$, use all 181 images (capping at available data).
    
    \item \textbf{Fine-tune}: Load joint-pretrained weights and continue training for 150 epochs with:
    \begin{itemize}
        \item \textbf{Learning rate}: $\text{lr}_0 = 0.0001$ (100× lower than pretraining, typical for fine-tuning)
        \item \textbf{Other params}: Same as pretraining (batch size, augmentation, etc.)
        \item \textbf{Validation}: Every epoch on full Kiwi validation set (60 images, fixed across all budgets)
        \item \textbf{Early stopping}: Patience = 20
    \end{itemize}
    
    \item \textbf{Evaluate}: Report mAP@50 (primary) and mAP@50:95, precision, recall, F1 on test set (121 images).
\end{enumerate}

\textbf{Rationale for parameters}: Lower learning rate (0.0001) prevents catastrophic forgetting of pre trained features. 150-epoch limit (vs. 300 for pretraining) is justified by convergence analysis: on small labeled sets, models plateau by epoch 100--120 (Section~\ref{sec:results}).

\subsection{Stage 3: Experimental Design and Baselines}

\subsubsection{Baselines}

We compare CrossPoll against three baselines:

\begin{table}[H]
\centering
\caption{Baseline Methods}
\label{tab:baselines}
\begin{tabular}{lll}
\toprule
\textbf{Baseline} & \textbf{Pretraining Source} & \textbf{Fine-tune Data} \\
\midrule
\textbf{Scratch} & Random init / YOLO architecture & Kiwi: $B$ images \\
\textbf{COCO Transfer} & COCO pretrained (yolov8n.pt) & Kiwi: $B$ images \\
\textbf{Single-Crop Transfer} & Pretrain on Apple only & Kiwi: $B$ images \\
\textbf{CrossPoll (Ours)} & Pretrain on Apple + Strawberry + Tomato & Kiwi: $B$ images \\
\bottomrule
\end{tabular}
\end{table}

Scratch and COCO Transfer are standard baselines. Single-Crop Transfer isolates the effect of multi-source pretraining: if multi-source is beneficial, CrossPoll should exceed single-crop pretraining.

\subsubsection{Experimental Protocol}

\begin{itemize}
    \item \textbf{Label budgets}: $B \in \{25, 50, 100, 200, 500\}$ (five points for sample-efficiency curve)
    \item \textbf{Random seeds}: Three seeds (42, 43, 44) for each method × budget combination
    \item \textbf{Total runs}: $4 \text{ methods} \times 5 \text{ budgets} \times 3 \text{ seeds} = 60$ total training runs
    \item \textbf{Metrics}: 
    \begin{itemize}
        \item Primary: mAP@50 (intersection-over-union threshold 0.5, standard in agricultural applications)
        \item Secondary: mAP@50:95, precision, recall, F1
        \item Efficiency: Training time (min), model size (MB)
    \end{itemize}
\end{itemize}

\textbf{Motivation for three seeds}: Sample efficiency experiments are sensitive to random initialization and stochastic subsampling. Three seeds provide confidence intervals and flag outliers (e.g., seed 43 exhibits high variance at $B=50$—a known phenomenon in few-shot learning).

\subsection{Implementation Details}

\subsubsection{Hardware and Software}

\begin{itemize}
    \item \textbf{Hardware}: Apple M1 Pro (8-core CPU, 16GB unified memory)
    \item \textbf{Software}: Python 3.14, PyTorch 2.11.0, Ultralytics 8.4.26, MPS backend
    \item \textbf{Optimization flags}: 
    \begin{itemize}
        \item \texttt{cache="ram"} (loads all training images into RAM for 20--40\% speedup)
        \item \texttt{cos\_lr=True} (cosine annealing, better convergence on small datasets)
        \item \texttt{workers=4} (optimal for macOS; avoids fork overhead with higher counts)
        \item \texttt{deterministic=False} (minor speedup, seeding with \texttt{seed} param still controls variance)
        \item \texttt{plots=False} (skip per-epoch plot generation)
    \end{itemize}
\end{itemize}

\subsubsection{Reproducibility}

All experiments use fixed random seeds (42, 43, 44). Weights are saved as PyTorch checkpoints. We compute metrics using Ultralytics' evaluation API, which implements standard COCO mAP computation. Configuration files and scripts are located in:

\begin{itemize}
    \item Config: \texttt{configs/experiments/crosspoll.yaml}
    \item Training script: \texttt{scripts/run\_baselines.py}
    \item Results: \texttt{results/summary.csv}
\end{itemize}

\subsection{Expected Performance and Validation}

Preliminary runs (11 of 30 completed as of writing) show:

\begin{itemize}
    \item \textbf{COCO Transfer (b100)}: mAP@50 = 0.91 (competitive baseline)
    \item \textbf{Scratch (b100)}: mAP@50 = 0.92 (strong on medium budgets, weaker at small budgets)
    \item \textbf{Scratch (b25)}: mAP@50 = 0.006 (severe catastrophic failure at tiny budgets)
\end{itemize}

These preliminary results support our hypothesis: at $B=100$, both COCO transfer and scratch reach similar performance (both converge to good models), but at tiny budgets ($B=25$), scratch fails catastrophically while transfer-based methods should succeed.
